\clearpage
\section*{Judge-response synthesis: a provenance-first AMP benchmark}
\subsection*{The primary finding is a dataset validity failure}
The label-provenance audit is the central result of this paper. In the archived UniProt scale-up, 21,742 of 24,628 AMP-labelled examples are unreviewed while all 24,628 sampled negatives are reviewed. Class and review status are coupled by construction. Sequence-only models can exploit sequence family, organism, annotation and collection biases correlated with that status even though no status field is passed as input. The frozen CNN's 0.921 test AUROC measures ranking within this confounded collection; it cannot be promoted to a clean AMP-recognition result. The APD6 0.926 uses independently sourced positives but the same reviewed UniProt negative population, so it does not independently establish cross-database generalization either. The post-hoc 239-reviewed-positive subgroup is selected and small; its score does not repair the training or comparator bias. The study preserves this source error and its measured consequences rather than burying them under more model tuning.

\subsection*{What a valid replacement benchmark must lock}
The next study must fix cohort definitions before fitting: reviewed KW-0929 positives; negatives matched on review status, organism strata and length, with explicit discussion that absence of KW-0929 is not confirmed non-AMP activity. Publish class-by-review-status and organism cross-tabs for every split. De-duplicate exact sequences, then define protein-family-disjoint train/validation/test partitions with an alignment-based cluster audit, preserving accessions, source queries, checksums and exclusion reasons. Split before tuning preprocessing, calibration or ensemble weights. Keep a fully untouched external cohort that is neither UniProt KW-0929 nor APD6, if one with compatible labels and assay definitions can be obtained. Report class balance, family overlap, confidence intervals respecting cluster dependence, low-FPR sensitivity and failures by biological family. These are proposed standards, not results achieved by the present paper. [PENDING: committed reviewed-only matched retrain, MMseqs2 audit and truly external result.]

\subsection*{Practitioner decision path}
\begin{enumerate}
\item If positives come from KW-0929, examine the exact query and accession review prefix before training. If class predicts review status by design, stop: an apparently high AUROC cannot isolate antimicrobial function.
\item If the negatives mean only ``not annotated AMP,'' say that. Match source, organism and length where possible, and test against experimentally characterized non-AMP controls. Do not call an unannotated protein inactive without an assay.
\item If near-identical proteins can cross splits, cluster by sequence alignment before partitioning and audit the remaining train/test homology. Exact ten-mer exclusion alone is a floor, not a family-level guarantee.
\item If a candidate is generated by maximizing the same model that scores it, treat that score as a search objective, not independent validation. Run a locked external predictor, a broader alignment novelty audit, toxicity checks and biological assays before claiming activity.
\end{enumerate}

\subsection*{Negative architecture and screen findings}
The tested GNN uses fixed chain-window edges; it does not incorporate three-dimensional contacts or evolutionary embeddings. Its poorer score than the CNN on tested regimes is an honest negative for this implementation, not evidence that structure-informed graphs are useless. A future contact-map or embedding experiment needs an independent predefined comparison; no such improvement is claimed here. The 852-sequence screen is post-hoc and enriched for secretion-like/cysteine-rich hits. Removing the signal peptide from one named sequence lowered its ensemble score from 0.9082 to 0.6299 in the committed ablation. That single perturbation is diagnostic, not a population-wide causal estimate. Original hits remain unmeasured hypotheses, and the pairwise 3-mer Jaccard filter is not a global novelty certificate. No MIC, hemolysis, potency or delivery assay was performed.

\subsection*{Unrun controls and comparator plan}
A per-family evaluation needs defensible annotations for anionic, disulfide-rich and other AMP classes before calculating a family score. Existing APD6 sensitivity 0.601 at 5\% false-positive rate cannot identify which family drives missed positives. The current discovery screen uses a simple average, not a calibrated posterior; validation-only weight fitting and Platt/isotonic calibration should be tested on a new matched cohort with an untouched test set. AMPlify and Macrel are relevant published tools, but this paper does not yet hold their predictions on identical examples under compatible model-training access. A same-row comparison with documented versions, allowed training data and low-FPR metrics is [PENDING]. Scrambled-sequence and experimentally non-AMP negative controls, physicochemical-channel ablation, explicit length-confound test, alignment-based database novelty checks, and DBAASP continuous MIC regression are also [PENDING]. A web screening interface would expose confidence and abstention, but is not counted as wet-lab validation. A locked future discovery protocol must specify reference database versions, candidate inclusion/exclusion rules, model weights, score threshold, toxicity triage and assay plan before viewing new scores; it cannot retroactively preregister these original candidates.

\subsection*{Twelve-slide judge narrative, editorial storyboard only}
This outline is not a rendered slide deck. Slide 1: question and claim boundary. 2: cohort/source diagram. 3: exact review-status cross-tab. 4: how the positive-query omission happened. 5: why the 0.921 and 0.926 measurements are confounded. 6: the small reviewed-only post-hoc check and why it does not fix training. 7: the negative GNN result. 8: exact-ten-mer leakage in the published Feb2020 folds. 9: hard-split primary 0.40 result and 0.35 sensitivity with mixed gates. 10: unvalidated candidate screen and signal-peptide ablation. 11: locked replacement-benchmark plan. 12: decision tree, limitations and testable next step. Numbers on slides must be reconciled to the latest committed results before presentation. [PENDING: designed, rendered and visually inspected 12-slide presentation.]
